% Generated by Listas/make_statements.py from pratica-lecture-06.md, item 23. Do not edit.
\textbf{Source.} Gibbons (1992), \emph{Game Theory for Applied Economists}, Ch. 2 Problems, Section 2.3, Problem 2.14, pp.~135--136 {[}PDF 142--143{]}. Reproduced verbatim (the book's spelling ``Rotemburg'' kept). Answer in the source: none (Gibbons prints no answers).



\begin{quote}

2.14. Suppose that demand fluctuates randomly in the infinitely repeated Bertrand game described in Problem 2.13: in each period, the demand intercept is \(a_H\) with probability \(\pi\) and \(a_L\) (\(< a_H\)) with probability \(1 - \pi\); demands in different periods are independent. Suppose that each period the level of demand is revealed to both firms before they choose their prices for that period. What are the monopoly price levels (\(p_H\) and \(p_L\)) for the two levels of demand? Solve for \(\delta^*\), the lowest value of \(\delta\) such that the firms can use trigger strategies to sustain these monopoly price levels (i.e., to play \(p_i\) when demand is \(a_i\), for \(i = H, L\)) in a subgame-perfect Nash equilibrium. For each value of \(\delta\) between 1/2 and \(\delta^*\), find the highest price \(p(\delta)\) such that the firms can use trigger strategies to sustain the price \(p(\delta)\) when demand is high and the price \(p_L\) when demand is low in a subgame-perfect Nash equilibrium. (See Rotemburg and Saloner 1986.)

\end{quote}

